%============================================================
% CHƯƠNG 10: ĐẠO LÀM CHA MẸ
%============================================================
\chapter{Đạo Làm Cha Mẹ (The Way of Parenthood)}

\begin{center}
  \textit{\color{truyenblue}``The most important thing a father can do for his children is to love their mother.''}\\
  \textit{(Điều quan trọng nhất một người cha có thể làm cho con cái là yêu thương mẹ của chúng.)}\\[4pt]
  \small\color{truyengold}--- Theodore Hesburgh ---
\end{center}
\ngancach

\section{Đôi Giày Chật}
\begin{truyen}{Đôi Giày Chật}{Hi Sinh}
\chuhoa{N}{ }gười cha mua cho con trai đôi giày mới dịp khai giảng, trong khi đôi giày của mình đã rách đế từ tháng trước.\\
\textit{(The father bought his son new shoes for the start of school, while his own shoes had worn-through soles since the month before.)}

``Bố không cần giày mới à?'' cậu bé hỏi.\\
\textit{(``Don't you need new shoes, Dad?'' the boy asked.)}

``Chân bố quen rồi,'' ông cười. ``Hơn nữa, \textbf{bố đi xe máy. Con đi bộ đến trường}. Con cần hơn.''\\
\textit{(``My feet are used to it,'' he smiled. ``Besides, \textbf{I ride a motorbike. You walk to school}. You need them more.'')}

Cậu bé không hiểu lúc đó. Mười năm sau, khi tự mua được giày cho bố, cậu nhớ lại và \textbf{khóc lần đầu tiên vì lý do hoàn toàn không phải buồn}.\\
\textit{(The boy did not understand then. Ten years later, when he could finally buy shoes for his father, he remembered and \textbf{cried for the first time for a reason that had nothing to do with sadness}.)}
\end{truyen}
\begin{baihoc}
  Sự hi sinh của cha mẹ thường không nói thành lời và chỉ được nhận ra khi đã đủ lớn để hiểu. Trân trọng những điều nhỏ bé đó trước khi quá muộn là điều con cái nên học.
\end{baihoc}

\section{Bản Thảo Cuốn Sách}
\begin{truyen}{Bản Thảo Cuốn Sách}{Ước Mơ}
\chuhoa{N}{ }gười mẹ viết lách từ hồi trẻ nhưng bỏ dở khi sinh con. Bản thảo nằm trong ngăn kéo hai mươi năm.\\
\textit{(The mother had written since her youth but stopped when she had children. The manuscript sat in a drawer for twenty years.)}

Khi con gái lớn phát hiện ra, cô đọc ngấu nghiến và hỏi: ``Sao mẹ không in thành sách?''\\
\textit{(When her grown daughter discovered it, she read it hungrily and asked: ``Why didn't you publish this?'')}

``Lúc đó bận nuôi con.''\\
\textit{(``I was busy raising you then.'')}

Con gái đưa bản thảo đến nhà xuất bản. Sáu tháng sau, sách ra đời với tên tác giả và lời đề tặng: \textbf{\textit{``Gửi con gái, người đã tìm lại giấc mơ của mẹ hộ mẹ.''}}\,\\
\textit{(The daughter took the manuscript to a publisher. Six months later, the book was released with the author's name and a dedication: \textbf{\textit{``To my daughter, who found my dream for me.''}}\,)}
\end{truyen}
\begin{baihoc}
  Cha mẹ thường gác lại ước mơ của mình để nuôi dưỡng ước mơ của con. Đôi khi món quà đẹp nhất con cái có thể tặng cha mẹ là giúp họ tìm lại những gì đã bỏ lại.
\end{baihoc}

\section{Cuộc Điện Thoại Đêm Khuya}
\begin{truyen}{Cuộc Điện Thoại Đêm Khuya}{Luôn Ở Đó}
\chuhoa{Đ}{ }ứa con gọi điện lúc hai giờ sáng, khóc mà không nói được lý do.\\
\textit{(The child called at two in the morning, crying without being able to explain why.)}

Người mẹ nghe máy ngay trong tiếng chuông đầu. Không hỏi nhiều. Chỉ nói: ``Mẹ ở đây rồi.''\\
\textit{(The mother answered on the first ring. Without many questions. Just said: ``I am here.'')}

Họ ngồi im ở hai đầu dây nửa tiếng.\\
\textit{(They sat in silence on either end of the line for half an hour.)}

Rồi con nói: ``Con ổn rồi mẹ. Cảm ơn mẹ.''\\
\textit{(Then the child said: ``I am okay now, Mom. Thank you.'')}

Người mẹ không ngủ lại được nhưng \textbf{nằm mỉm cười trong bóng tối}.\\
\textit{(The mother could not fall back asleep but \textbf{lay smiling in the darkness}.)}
\end{truyen}
\begin{baihoc}
  Cha mẹ không cần hiểu mọi nỗi đau của con để có thể an ủi con. Đôi khi chỉ cần hiện diện, không phán xét, không hỏi han là điều con cái cần nhất giữa đêm khuya.
\end{baihoc}

\section{Buổi Tập Võ Cuối Cùng}
\begin{truyen}{Buổi Tập Võ Cuối Cùng}{Buông Tay}
\chuhoa{T}{ }rong tám năm, người cha đưa con trai đi tập võ mỗi sáng thứ Bảy. Rồi cậu bé nói muốn nghỉ để theo học nhạc.\\
\textit{(For eight years, the father drove his son to morning martial arts practice every Saturday. Then the boy said he wanted to quit to study music instead.)}

Nhiều bậc cha mẹ sẽ tranh luận về sự kiên trì. Người cha này im lặng một hôm.\\
\textit{(Many parents would argue about perseverance. This father was quiet for a day.)}

``Con chắc chắn chưa?''\\
\textit{(``Are you sure?'')}

``Dạ chắc.''\\
\textit{(``Yes, I am sure.'')}

``Vậy bố mua đàn cho con cuối tuần này.''\\
\textit{(``Then I will buy you an instrument this weekend.'')}

Mười lăm năm sau, người con biểu diễn trong dàn nhạc giao hưởng quốc gia. \textbf{Cha ngồi hàng đầu, khóc không che giấu}.\\
\textit{(Fifteen years later, the son performed with the national symphony orchestra. \textbf{His father sat in the front row, crying openly}.)}
\end{truyen}
\begin{baihoc}
  Yêu con không phải là bắt con đi theo con đường mình đã vạch ra. Biết trao quyền quyết định cho con đúng lúc là một trong những hành động yêu thương trưởng thành nhất của cha mẹ.
\end{baihoc}

\section{Người Cha Không Biết Chữ}
\begin{truyen}{Người Cha Không Biết Chữ}{Kiên Cường}
\chuhoa{Ô}{ }ng không biết đọc, không biết viết. Nhưng ông thuộc lòng tên của cô giáo con, tên của từng môn học con học, tên của người bạn thân nhất của con.\\
\textit{(He could neither read nor write. But he knew by heart his child's teacher's name, every subject the child studied, the name of the child's closest friend.)}

Mỗi tối ông ngồi bên con đang làm bài, không giúp được gì về chữ nghĩa nhưng \textbf{pha trà, bẻ bánh, đuổi muỗi}.\\
\textit{(Every evening he sat beside his child doing homework, unable to help with any of the words but \textbf{making tea, breaking crackers, chasing mosquitoes}.)}

``Con học bài đi, có bố đây.''\\
\textit{(``Study now, your father is here.'')}

Năm con đậu đại học, ông đến trường nghe loa đọc tên con mà không hiểu chữ nào trên tấm bảng kết quả. Nhưng \textbf{ông nghe thấy tên con và cúi mặt xuống khóc ròng}.\\
\textit{(The year his child passed the university entrance exam, he came to the school and heard the loudspeaker read his child's name without understanding a single character on the results board. But \textbf{he heard his child's name and bowed his head and wept}.)}
\end{truyen}
\begin{baihoc}
  Học vấn không đo lường được tình yêu thương. Người cha không biết chữ nhưng hiểu rõ điều con cần hơn bất kỳ sách giáo khoa nào có thể dạy.
\end{baihoc}

\section{Hai Nửa Bánh Mì}
\begin{truyen}{Hai Nửa Bánh Mì}{Bình Đẳng}
\chuhoa{T}{ }rong nhà có hai con, một chiếc bánh mì duy nhất và không đủ tiền mua thêm.\\
\textit{(In the household there were two children, one single loaf of bread and no money to buy more.)}

Người mẹ chia đôi đều đặn rồi lấy phần vụn còn lại trên thớt cho mình.\\
\textit{(The mother divided it evenly, then took the crumbs left on the cutting board for herself.)}

Đứa lớn hỏi: ``Mẹ không đói à?''\\
\textit{(The older child asked: ``Aren't you hungry, Mom?'')}

``Mẹ ăn rồi.''\\
\textit{(``I already ate.'')}

Cả hai đứa trẻ lớn lên biết rằng \textbf{lời đó không bao giờ là thật}, nhưng cũng biết rằng đó là \textbf{cách mẹ nói ``tao yêu mày'' theo ngôn ngữ của người nghèo}.\\
\textit{(Both children grew up knowing \textbf{that statement was never true}, but also knowing it was \textbf{the way their mother said ``I love you'' in the language of the poor}.)}
\end{truyen}
\begin{baihoc}
  Tình yêu của cha mẹ đôi khi được nói bằng những lời dối trắng và những nhường nhịn im lặng. Đó là ngôn ngữ mà chỉ khi lớn lên ta mới học đọc được.
\end{baihoc}

\section{Tin Nhắn Không Gửi}
\begin{truyen}{Tin Nhắn Không Gửi}{Tình Yêu Thầm Lặng}
\chuhoa{S}{ }au khi mẹ mất, người con dọn điện thoại cũ của bà và tìm thấy hàng chục tin nhắn đã soạn nhưng không bao giờ được gửi đi.\\
\textit{(After the mother passed away, the child cleaned her old phone and found dozens of messages drafted but never sent.)}

\textit{``Con ơi, mẹ nhớ con lắm. Nhưng mẹ sợ làm phiền con.''}\\
\textit{(``My child, I miss you so much. But I was afraid of bothering you.'')}

\textit{``Hôm nay mẹ nấu món con thích, ước gì con ở đây.''}\\
\textit{(``Today I cooked your favourite dish and wished you were here.'')}

\textit{``Con có ổn không? Mẹ muốn hỏi mà sợ con bận.''}\\
\textit{(``Are you okay? I wanted to ask but was afraid you were busy.'')}

Người con đọc từng tin nhắn, \textbf{mỗi chữ như một mũi kim}. Và tự hỏi tại sao mình không gọi về nhiều hơn khi còn kịp.\\
\textit{(The child read each message, \textbf{each word like a needle}. And wondered why they had not called home more often while there was still time.)}
\end{truyen}
\begin{baihoc}
  Cha mẹ thường sợ làm phiền con cái hơn là sợ cô đơn. Đừng để họ phải soạn những tin nhắn không bao giờ dám gửi. Hãy gọi về trước khi không còn ai bắt máy nữa.
\end{baihoc}

\section{Lớp Học Của Bố}
\begin{truyen}{Lớp Học Của Bố}{Kiên Nhẫn}
\chuhoa{K}{ }hi con gái bắt đầu dạy bố dùng điện thoại thông minh, cô mất kiên nhẫn chỉ sau mười phút.\\
\textit{(When the daughter started teaching her father to use a smartphone, she lost patience after just ten minutes.)}

``Bố hỏi mãi một câu.''\\
\textit{(``You keep asking the same question.'')}

Ông im lặng, nhìn xuống rồi nói nhẹ: ``Hồi con mới học đi, con té mấy trăm lần. Bố đếm từng cái. Bố không bao giờ thở dài.''\\
\textit{(He was quiet, looked down, then said softly: ``When you were first learning to walk, you fell hundreds of times. Father counted every one. Father never sighed.'')}

Cô gái xấu hổ, ngồi lại, \textbf{dạy từ đầu với giọng nhẹ hơn}.\\
\textit{(The daughter was ashamed, sat back down, and \textbf{taught from the beginning with a gentler voice}.)}
\end{truyen}
\begin{baihoc}
  Cha mẹ từng kiên nhẫn vô tận khi ta còn nhỏ. Đến lúc họ cần ta kiên nhẫn với họ, đó là cơ hội để trả một phần món nợ ân tình không bao giờ trả hết được.
\end{baihoc}

\section{Bức Tranh Dán Tường}
\begin{truyen}{Bức Tranh Dán Tường}{Tự Hào}
\chuhoa{T}{ }rong phòng khách nhà một người thợ hồ nghèo, bức tranh đẹp nhất trên tường không phải ảnh phong cảnh hay danh nhân.\\
\textit{(In the living room of a poor bricklayer, the most prominent thing on the wall was not a landscape or a portrait of a famous figure.)}

Đó là tờ giấy khen học sinh giỏi của con gái, đã ố vàng và nhàu nát nhưng được bọc nhựa ép kỹ lưỡng.\\
\textit{(It was his daughter's academic merit certificate, yellowed and creased but carefully laminated.)}

Khách đến nhà lần nào ông cũng chỉ vào: ``Con gái tôi đấy.''\\
\textit{(Every time visitors came, he pointed to it: ``That is my daughter.'')}

Con gái sau này trở thành kỹ sư, treo bằng đại học lên tường riêng. Nhưng \textbf{bức tường bố treo mới là bức tường con trân trọng nhất}.\\
\textit{(His daughter later became an engineer, hanging her university degree on her own wall. But \textbf{her father's wall was the one she treasured most}.)}
\end{truyen}
\begin{baihoc}
  Niềm tự hào của cha mẹ về con cái không đo bằng bằng cấp hay địa vị, mà bằng tình yêu thuần túy dành cho từng bước trưởng thành nhỏ nhoi của con. Đó là điều không một thành tích nào có thể thay thế.
\end{baihoc}

\section{Người Mẹ Sau Cánh Gà}
\begin{truyen}{Người Mẹ Sau Cánh Gà}{Tự Hào Thầm Lặng}
\chuhoa{Đ}{ }êm diễn kịch của trường, cô bé vào vai phụ: chỉ một câu thoại và ba giây trên sân khấu.\\
\textit{(On the school play night, the girl had a supporting role: just one line and three seconds on stage.)}

Giữa hàng trăm khán giả, người mẹ đứng dậy vỗ tay khi con bước ra, \textbf{dù con chỉ đứng ở góc bên phải màn hình}.\\
\textit{(Among hundreds of audience members, the mother stood up and clapped when her daughter stepped out, \textbf{even though the girl stood only at the far right corner of the stage}.)}

Bạn bè cô bé thì thầm: ``Mẹ bạn vỗ tay to nhất.''\\
\textit{(The girl's friends whispered: ``Your mom clapped the loudest.'')}

Cô bé nhìn xuống, tìm ánh mắt mẹ và \textbf{lần đầu tiên không còn thấy vai mình nhỏ nữa}.\\
\textit{(The girl looked down into the audience, found her mother's eyes and \textbf{for the first time no longer felt her role was small}.)}
\end{truyen}
\begin{baihoc}
  Cha mẹ thấy con ở những chỗ người khác không nhìn thấy và vỗ tay cho những khoảnh khắc người khác bỏ qua. Đó là tình yêu bền vững nhất trên đời.
\end{baihoc}
